雅化的竞争位置由两条不同的链条决定：锂盐业务需要把矿源、转化能力和客户资格转成稳定吨利润；民爆业务需要把许可、区域网络和项目执行转成可回收的服务现金。两条链均有进入壁垒，但壁垒本身不自动构成 2026 年收入、成本或客户份额。本章将可验证的经营事实与不应提前计入的叙事分开。

\section{锂盐：供应路径与客户资格是必要条件，不是吨成本证明}

公司披露拥有李家沟与 KMC 等资源布局、多个锂盐生产基地，以及面向车企、电池和正极企业的客户渠道。李家沟优先供应权和资源储量说明未来供应保障可能改善，客户名单说明产品已进入较高门槛的供应体系；但截至本报告截止日，雅化未披露 2026 年来自各资源的实际供给量、结算价格、单位成本、客户分配或剩余订单。因此，模型只采用历史锂盐产品篮子，不采用“资源自给更高”或“客户放量”作为基准超额利润。

\begin{exhibitbox}[表 6-1：雅化锂盐链条的估值权限]
\begin{tabularx}{\exhibitboxwidth}{L{2.9cm}X X}
\toprule
链条环节 & 已确认的经营含义 & 基准模型的处理 \\
\midrule
资源与优先供应 & 李家沟权益/优先供应权、KMC 等资源布局构成潜在供应路径 & 无实际供给、成本和自供比例时，不给成本优势信用 \\
锂盐转化 & 三个基地、近 13 万吨设计综合产能说明制造能力边界 & 设计产能不替代实际销量或利用率 \\
客户资格 & 面向车企、电池和正极材料客户的长期渠道存在 & 无客户份额、价格和订单时，不给收入增量信用 \\
价格与套保 & 公司披露参考市场定价并使用期货对冲价格波动 & 无实际 ASP、套保损益和库存数据时，不用作利润预测锚 \\
\bottomrule
\end{tabularx}
\end{exhibitbox}
\sourcenote{公司《2025 年年度报告》业务模式、资源与客户披露；客户、资源和套保的收入/成本贡献均按未披露处理。}

\section{民爆：区域牌照与一体化能力提供防线，但需穿透项目现金}

民爆产品受许可与运输半径约束，雅化在四川、内蒙古、山西、西藏、新疆等区域布局，并将产品销售与爆破矿服结合。这一网络比单纯产能更有价值，因为产品与服务可以共同覆盖矿山和基建项目全周期；但新增项目金额、回款节奏、服务毛利改善和海外扩张均没有形成可审计的 2026 年增量证据。

\begin{keyinsight}[竞争结论]
雅化的相对优势应表述为“具备从锂矿/锂盐到高门槛客户、以及从民爆产品到矿服的经营平台”，而不是“已经获得更低成本或更大订单”。前者可以解释为何公司值得被持续覆盖，后者只有在正式披露的量、价、成本和现金数据出现后才能提高估值。
\end{keyinsight}

\section{同业比较：低估值并不自动代表更高质量}

天齐锂业、赣锋锂业和盛新锂能均受锂价与行业库存周期影响，但资源权益、产品结构、海外资产、现金流和历史盈利波动不同。以 2026E 聚合 PE 观察，三者为 11.71 倍、13.51 倍和 13.16 倍，均值 12.79 倍、中位数 13.16 倍。它们只能作为估值倍数的低权重校准：雅化的民爆防御性利润可降低部分周期波动，但负经营现金流与未验证的资源成本也应限制其估值溢价。

\begin{exhibitbox}[表 6-2：可比框架——比较的是盈利质量，不是名称]
\begin{tabularx}{\exhibitboxwidth}{L{2.6cm}L{2.5cm}X X}
\toprule
对象 & 主要经营特征 & 可用于雅化估值的启发 & 不可直接迁移的部分 \\
\midrule
天齐锂业 & 资源与锂盐一体化、海外权益影响显著 & 大型锂企的周期估值区间 & 资产结构、权益法与资源成本 \\
赣锋锂业 & 资源布局和多元锂产品并存 & 价格周期与资源保障的折现方式 & 产品结构、海外项目和现金流 \\
盛新锂能 & 锂盐周期弹性更集中 & 锂价反弹下的利润敏感性 & 资源、客户和产能利用率 \\
雅化集团 & 锂盐加民爆/矿服双主业 & 民爆可形成历史利润底座 & 未验证资源自供、客户份额和项目回款 \\
\bottomrule
\end{tabularx}
\end{exhibitbox}
\sourcenote{同行 2026E 聚合估值快照；公司公开年报与行业资料。同行 PE 为低权重比较，不构成雅化的直接目标价。}
